\subsection{Topological vector spaces over a field with a valuation}

Throughout let K be a (not necessarily commutative) field, v a valuation on K and G its order group. K is given the topology \(T_{v}\) .

PROPOSITION 3. Let E be a left topological vector space over K which is Hausdorff and of dimension 1. Suppose that v is not improper. For all \(x_{0} \neq 0\) in E, the mapping \(a \mapsto ax_{0}\) of \(K_{s}\) onto E is a topological isomorphism.

This mapping is a continuous algebraic isomorphism. It is sufficient to show that it is bicontinuous. Let \(\alpha \in \mathbf{G}\) . We need to show that there exists a neighbourhood V of 0 in E such that the relation \(ax_0 \in \mathbf{V}\) implies \(v(a) > \alpha\) . Let \(a_0 \in \mathbf{K}^*\) be such that \(v(a_0) = \alpha\) . As E is Hausdorff, there exists a neighbourhood W of 0 in E such that \(a_0 x_0 \notin \mathbf{W}\) . As \(v\) is not improper, there exist a neighbourhood W' of 0 in E and an element \(\beta\) of G such that the relations \(y \in \mathbf{W}'\) , \(v(a) \geqslant \beta\) imply \(ay \in \mathbf{W}\) . Let \(a_1 \in \mathbf{K}^*\) be such that \(v(a_1) = -\beta\) . The relations \(ax_0 \in a_1^{-1} \mathbf{W}'\) and \(v(a) \leqslant \alpha\) imply \(a_1 ax_0 \in \mathbf{W}'\) and \(v(a_0 a^{-1} a_1^{-1}) = \alpha + \beta - v(a) \geqslant \beta\) and hence \(a_0 x_0 = a_0 a^{-1} a_1^{-1}(a_1 ax_0) \in \mathbf{W}\) , which is absurd; in other words, the relation \(ax_0 \in a_1^{-1} \mathbf{W}\) implies \(v(a) > \alpha\) .

COROLLARY. Let E be a left topological vector space over K, H a closed hyperplane of E and D a 1-dimensional vector subspace of E an algebraic supplement of H. Suppose that v is not improper. Then D is a topological supplement of H.

Taking account of Propositions 1 and 3, the proof is the same as that of Topological Vector Spaces, Chapter I, § 2, Corollary 2 to Theorem 1.

PROPOSITION 4. Suppose that \(v\) is not improper and \(K\) is complete. Let \(E\) be a left topological vector space over \(K\) , which is Hausdorff and of finite dimension \(n\) . For every basis \((e_i)_{1 \leqslant i \leqslant n}\) of \(E\) over \(K\) , the mapping \((a_i) \mapsto \sum_{i=1}^{n} a_i e_i\) of \(K_s^n\) onto \(E\) is a topological vector space isomorphism.

Taking account of Proposition 3 and its corollary, the proof is the same as that of Topological Vector Spaces, Chapter I, § 2, Theorem 2.

COROLLARY. Suppose that v is not improper and K is complete. Let E be a Hausdorff topological vector space over K and F a finite-dimensional vector subspace of E. Then F is closed.

F is complete.

